\documentclass[aspectratio=169]{beamer}
\usetheme{metropolis}
\usepackage{nexus_theme}
\usepackage{listings}
\usepackage{booktabs}

\title{Critical Reading \& Synthesis}
\subtitle{Module 3.1: Analyzing Literature}
\author{Nexus Scholar Suite --- Modern Research Methodologies}
\date{}

\begin{document}

% --- TITLE SLIDE ---
\begin{frame}[plain]
  \titlepage
\end{frame}

% --- LEARNING OBJECTIVES ---
\begin{frame}{What You'll Learn}
  By the end of this lesson, you will be able to:
  \begin{enumerate}
    \item \textbf{Apply} an extraction template for comparable data.
    \item \textbf{Build} a synthesis matrix (concept matrix).
    \item \textbf{Distinguish} summary from synthesis.
    \item \textbf{Identify} convergence, divergence, gaps, and evolution.
  \end{enumerate}
\end{frame}

% --- THE PROBLEM ---
\begin{frame}{Summary $\neq$ Synthesis}
  \begin{columns}[T]
    \begin{column}{0.48\textwidth}
      \begin{alertblock}{Summary (Don't do this)}
        ``Smith (2020) found X.\\
        Jones (2021) found Y.\\
        Lee (2022) found Z.''
        \vspace{0.2cm}
        
        \textit{A grocery list of individual findings.}
      \end{alertblock}
    \end{column}
    \begin{column}{0.48\textwidth}
      \begin{exampleblock}{Synthesis (Do this)}
        ``Three studies converge on X (Smith 2020; Jones 2021; Lee 2022), though Patel (2023) challenges this, arguing Y.''
        \vspace{0.2cm}
        
        \textit{A conversation between studies.}
      \end{exampleblock}
    \end{column}
  \end{columns}
\end{frame}

% --- EXTRACTION TEMPLATE ---
\begin{frame}{Step 1: The Extraction Template}
  \textbf{Problem:} Free-form notes on 50 papers = 50 incomparable formats.\\
  \textbf{Solution:} A standardized table for every paper.
  \vspace{0.3cm}
  
  \begin{table}
  \footnotesize
  \begin{tabular}{lp{8cm}}
  \toprule
  \textbf{Column} & \textbf{What to Extract} \\
  \midrule
  Author / Year & Citation reference \\
  Research Question & What did they ask? \\
  Methodology & Quantitative, qualitative, mixed? \\
  Sample / Data & Who or what did they study? \\
  Key Findings & The core result \\
  Limitations & What they acknowledged or missed \\
  Relevance & How does this connect to MY thesis? \\
  \bottomrule
  \end{tabular}
  \end{table}
\end{frame}

% --- SYNTHESIS MATRIX ---
\begin{frame}{Step 2: The Synthesis Matrix}
  \textbf{Flip the table:} Rows = themes, Columns = papers.
  \vspace{0.3cm}
  
  \begin{table}
  \footnotesize
  \begin{tabular}{l|c|c|c|c}
  \toprule
  \textbf{Theme} & \textbf{Smith '20} & \textbf{Jones '21} & \textbf{Lee '22} & \textbf{Patel '23} \\
  \midrule
  Data quality & Addressed & Addressed & --- & Addressed \\
  Model architecture & CNN & RNN & Transformer & CNN \\
  Ethics / bias & --- & --- & --- & Addressed \\
  Sample size & 500 & 1200 & 300 & 8000 \\
  \bottomrule
  \end{tabular}
  \end{table}
  \vspace{0.2cm}
  \textit{Empty cells (---) are not failures. They are gaps --- and gaps are where your thesis lives.}
\end{frame}

% --- FOUR MOVES ---
\begin{frame}{The Four Synthesis Moves}
  When writing from the matrix, use these intellectual moves:
  \vspace{0.3cm}
  \begin{enumerate}
    \item \textbf{Convergence:} ``Three studies independently found that...''
    \item \textbf{Divergence:} ``However, Patel (2023) challenges this...''
    \item \textbf{Gap:} ``No study has addressed the impact of X on Y.''
    \item \textbf{Evolution:} ``Early work (2015--2018) focused on X, but the field has shifted toward Y.''
  \end{enumerate}
\end{frame}

% --- EXERCISE PREVIEW ---
\begin{frame}{Your Turn: Exercise}
  \begin{exampleblock}{Exercise: The Synthesis Matrix Builder}
    Read 5 provided abstracts, build an extraction template, identify themes, and construct a synthesis matrix. Then write one synthesis paragraph using all four moves.
  \end{exampleblock}
\end{frame}

% --- STANDOUT CLOSE ---
\begin{frame}[standout]
  Your literature review is not a grocery list.\\
  It's a conversation --- and you need to tell\\
  your reader what the room is saying.
\end{frame}

\end{document}
